\section{Discussion and Conclusion}
\label{sec:discussion}

\paragraph{What the theory licenses.}
Our fixed-partition results (Theorems~\ref{thm:ratio}--\ref{thm:sparsification}) are correct and verified to high precision, and they should be read for exactly what they state: degree-biased sampling shifts edge composition and null-model mass in a direction determined by the score gap $\delta$, for whatever partition $\delta$ is measured against.
The configuration-model controls show that this premise, and therefore the conclusion, holds as readily on graphs without communities as on real networks.
The theory is a quantitative account of degree mechanics; treating it as a theory of community clarification was an over-interpretation---one we made ourselves, and one the literature on preprocessing-for-clustering is structurally exposed to, since the certifying statistics look meaningful and the standard evaluations all move in the flattering direction.

\paragraph{What we recommend.}
For practitioners: degree-based sparsification remains a legitimate \emph{scaling} tool at mild retention, where the fixed-objective cost is small (Section~\ref{sec:boundary-retention}), and seeded refinement is worth trying when clustering budgets allow a second pass---but no quality improvement should be assumed, and none should be claimed without the controls of Section~\ref{sec:protocol}.
For evaluators and reviewers: the three controls (fixed objective, granularity matching, compute matching) plus the configuration-model null are each a few lines of code, and each reversed at least one published-style conclusion in this paper.
For benchmark designers: standard LFR is doubly unrepresentative for sparsification studies---nearly homogeneous in degree and uncoupled between degree and community membership---and controlled degree--community coupling of the kind used here (or in richer generators) is required for synthetic conclusions to transfer.

\paragraph{Limitations.}
Our scope is modularity-based detection with Leiden under DSpar-family sparsification; other objectives (description-length, flow-based), other detectors, and similarity- or resistance-based sparsifiers may behave differently, and our protocol is designed to test exactly that.
The boundary characterization rests on 15 networks: sufficient to establish existence (2 genuine successes) and to defeat the tested predictors, insufficient to establish what does predict success.
Reference-label experiments inherit the known gap between metadata and structural communities~\cite{peel2017ground, hric2014community}.

\paragraph{Open problems.}
Why \texttt{email-Enron} and \texttt{com-Youtube}?
Both gains are real, robust, and---by every statistic we computed---unforecastable; understanding what landscape property lets hub-edge removal unlock better optima on precisely these graphs is the concrete open question this work leaves.
A second is constructive: whether a sparsifier can be \emph{designed} so that its selection signal separates real networks from their configuration models, rather than inheriting the degree sequence's behavior.

\paragraph{Conclusion.}
We set out to determine when sparsification helps community detection.
The answer has three parts.
The mechanism by which degree-based sparsification reshapes graphs is real, provable, and controllable.
The standard evidence that this mechanism improves detection is artifactual---attributable to cross-graph scoring, granularity, and unmatched computation, and reproduced in full by configuration models without community structure.
And a genuine benefit exists but is rare and currently unpredictable: two networks in fifteen, small margins, surviving every control we could construct.
Sparsification, in short, is a scaling tool that occasionally---and so far inexplicably---doubles as an optimizer's heuristic; it is not a general clarifier of community structure.
The evaluation protocol that separates these outcomes costs little, and we hope it becomes routine.

\section*{Acknowledgment}

This research was funded in part by NSF grant numbers CCF-2109988, OAC-2402560, and CCF-2453324.

\appendix

\section*{Reproducibility}
All experiments in this paper are reproducible from the accompanying repository: per-experiment scripts, per-run CSV outputs, fixed seeds, and the calibrated sampler implementation, together with the audit artifacts (including bit-level verification of our reimplementation of the with-replacement sampler against the reference implementation).
Realized retention is logged for every sparsification call.
Dataset sources: SNAP collections for the large networks and standard repositories for the classical labeled graphs, as cited per table.
